\chapter*{Introduction} \label{ch:cap1}







\section{A worthy note on context}

\begin{center}
    \centering
    \textit{The miracle of the appropriateness of the language of mathematics for the formulation of the laws of physics is a gift which we neither understand nor deserve.}
\end{center}

Eugene Wigner expressed his admiration in his 1960 article in \textit{Communications in Pure and Applied Mathematics.} It is a sentiment many would agree with: the effectiveness of mathematics in describing Nature is both remarkable and fascinating. Moreover, over the last century its branches have reached into areas where, not so long ago, mathematicians were considered too abstract to be of much use in business terms.

In 1827, Robert Brown, while studying the behaviour of pollen grains suspended in water, noticed something unexpected: the grains wandered along irregular, hard-to-describe paths. Physics was not prepared to tackle such a problem yet. What was needed belonged to what we now call \textit{statistical mechanics}, the study of collective behaviour of particles and molecules, intimately tied to \textit{thermodynamics.} The decisive step came later, in 1905, when Albert Einstein provided a theory of this motion — the same year he also explained the photoelectric effect, introduced the Special Relativity theory and provided the first appearance of the mass-energy equivalence formula, $E=mc^2$. What a year.

Regarding the polen grains, a parallel story took place in finance. Louis Bachelier, a doctoral student of Henri Poincaré, was interested in modelling stock prices. His model, which we now recognise as essentially Gaussian, worked with some shaky concepts on probability theory (Kolmogorov and his peers were not born yet, even Lebesgue had not published his work on measure theory yet). Worse, it allowed negative prices — a point critics did not miss.

The common thread in both stories is \emph{Brownian motion}. This stochastic process is a central object in probability: geometers find its fractal dimension fascinating, while probabilists value the rich family of processes built upon it. The idea of treating “random noise” as a genuine mathematical object took shape gradually, but it was Kiyoshi Itô’s work on \emph{stochastic integration} that truly changed the picture.

Finance joined somewhat later, but with undeniable impact. Today mathematics is everywhere in banks, consultancy and audit firms, underpinning the valuation of derivatives, that is, assigning a \emph{fair price} (according to market data) to contracts involving random future cash flows. This pricing orientation  will be the focus of the present thesis, among many other applications in finance, such as risk management and portfolio optimization.









\section{Objective and Scope}


The primary objective of this Thesis is to provide a solid review of \textit{Volatility}, aiming to fill the gap between those who are strictly Mathematics-oriented—who often lack expertise in Finance—and those who are Finance-oriented and do not fully appreciate the depth of the ideas they are working with. For this reconciliatory purpose, the Thesis embraces a clear and rigorous mathematical mindset, while maintaining continuous references to practical and financial interpretations. 




\hyperref[part:part1]{Part I} aims to provide a thorough introduction to the subject. The core mathematical ideas, rooted in Probability Theory and forming the basis of modern Financial Theory and basic Pricing Theory, are presented in Chapter \ref{ch:cap2}. Particular attention is paid there to the regularity of Brownian trajectories, since it is precisely the delicate balance between roughness and tractability --- non-vanishing quadratic variation on the one hand, unbounded total variation on the other --- that both makes Itô's theory necessary and marks the limits of its reach.

\hyperref[part:part2]{Part II} is devoted to more recent developments in volatility modelling, including Rough Paths and Rough Volatility theories (Chapter \ref{ch:cap5}). The latter is introduced from a general and mathematical perspective: although it finds direct applications in Finance, the theory of rough paths has far-reaching implications in other areas of mathematics as well. These chapters aim to highlight the necessity of more sophisticated tools to describe and operate in the market.

Finally, these theoretical constructs are brought to practice in Chapter \ref{ch:cap6}, where a calibration of \hl{XXX} is performed. Here, the models are numerically implemented and compared, providing an applied perspective and illustrating the strengths and challenges associated with each approach. Some final conclusions are left for Chapter \ref{ch:cap7}, where a general overview of the work done will be given, along with reflections on possible directions for further research.

\newpage
%\rule{0.95\textwidth}{1.5pt}

